\documentclass[11pt]{article}
\ifdefined\pdftrailerid\pdftrailerid{}\fi
\usepackage[T1]{fontenc}
\usepackage[utf8]{inputenc}
\usepackage{lmodern}
\usepackage[margin=1in]{geometry}
\usepackage{amsmath,amssymb,amsthm,mathtools}
\usepackage{enumitem,booktabs,array,longtable,needspace}
\usepackage{xcolor,microtype,etoolbox,xurl}
\usepackage[colorlinks=true,linkcolor=blue!45!black,citecolor=green!35!black,urlcolor=blue!55!black]{hyperref}
\usepackage[capitalize,nameinlink]{cleveref}
\hypersetup{pdftitle={Frozen Gap Areas and the Sharpness of Source-Cover Certificates},pdfauthor={Hans Brende and Claude Opus 5.5 (Anthropic)},pdfsubject={Frozen gap-area certificates for linear-time tape transposition, October 2026},pdfkeywords={matrix transposition, multitape Turing machines, source covers, circular gap area, Lean}}
\numberwithin{equation}{section}
% ---------------------------------------------------------------------------
% Theorem environments
% ---------------------------------------------------------------------------
\theoremstyle{plain}
\newtheorem{theorem}{Theorem}[section]
\newtheorem{lemma}[theorem]{Lemma}
\newtheorem{corollary}[theorem]{Corollary}
\newtheorem{proposition}[theorem]{Proposition}
\newtheorem{conjecture}[theorem]{Conjecture}
\newtheorem{question}[theorem]{Question}
\newtheorem{problem}[theorem]{Problem}
\theoremstyle{definition}
\newtheorem{definition}[theorem]{Definition}
\newtheorem{remark}[theorem]{Remark}
\newtheorem{example}[theorem]{Example}

% Formalized results use outlines; unformalized results retain complete written proofs.
\newenvironment{outline}[1][Outline of the formal proof]{\begin{proof}[#1]}{\end{proof}}
% Results that are not formalized get complete proofs, labelled as such.
\newenvironment{informalproof}{\begin{proof}[Proof (not formalized)]}{\end{proof}}

% ---------------------------------------------------------------------------
% Lean references
%   \lean{name}            a declaration cited with its source paper or namespace
%   \formal{file}{names}   a relative source file and comma-separated declaration names
%   \partformal{file}{names}{gap}
%                          a numbered result whose formal counterpart covers all but the
%                          part described in {gap}; the gap is proved or cited in the text
%   \notformalized         a numbered result that has no formal counterpart
%   \cited{where}          a result of an earlier paper used without proof: where it is proved,
%                          and whether it is Lean-verified there
% ---------------------------------------------------------------------------
% Lean names are typeset like paths, with line breaks allowed after "_", "." and "/".
\DeclareUrlCommand\leanname{\def\UrlFont{\ttfamily\upshape}}
\renewcommand{\UrlBreaks}{\do\_\do\.\do\/\do\-}
\newcommand{\lean}[1]{\leanname{#1}}
\newcommand{\leanfile}[1]{\leanname{#1}}
\newcommand{\leanlist}[1]{\begingroup\def\leansep{}%
  \renewcommand*{\do}[1]{\leansep\leanname{##1}\def\leansep{, }}\docsvlist{#1}\endgroup}
\newcommand{\formal}[2]{\par\noindent{\small\upshape\raggedright\textsf{Lean:} \leanlist{#2}\unskip
  \nobreak\hfil\penalty50\hskip1em\hbox{}\nobreak\hfill[\leanfile{#1}]\parfillskip=0pt\par}%
  \smallskip\noindent\ignorespaces}
\newcommand{\partformal}[3]{\par\noindent{\small\upshape\raggedright\textsf{Lean (partial):} \leanlist{#2}\unskip
  \nobreak\hfil\penalty50\hskip1em\hbox{}\nobreak\hfill[\leanfile{#1}]\parfillskip=0pt\par}%
  \par\noindent{\small\upshape\textsf{Not formalized:} #3}\par\smallskip\noindent\ignorespaces}
\newcommand{\notformalized}{\par\noindent{\small\upshape\textsf{Not formalized;} a complete proof is given
  below.}\par\smallskip\noindent\ignorespaces}
\newcommand{\cited}[1]{\par\noindent{\small\upshape\textsf{Cited:} #1}\par\smallskip\noindent\ignorespaces}

% ---------------------------------------------------------------------------
% Notation
% ---------------------------------------------------------------------------
\newcommand{\N}{\mathbb{N}}
\newcommand{\Z}{\mathbb{Z}}
\newcommand{\Q}{\mathbb{Q}}
\newcommand{\R}{\mathbb{R}}
\newcommand{\C}{\mathbb{C}}
\newcommand{\F}{\mathbb{F}}
\newcommand{\T}{\mathsf{T}}
\newcommand{\clog}[1]{\lceil \log_2 #1 \rceil}
\newcommand{\abs}[1]{\lvert #1 \rvert}
\newcommand{\card}[1]{\lvert #1 \rvert}
\newcommand{\eps}{\varepsilon}
\newcommand{\prev}{\mathrm{prev}}
\newcommand{\cost}{\mathrm{cost}}
\newcommand{\gn}{\mathrm{gn}}
\newcommand{\Sh}{\mathrm{Sh}}
\newcommand{\cl}{\mathrm{cl}}
\newcommand{\Aut}{\mathrm{Aut}}
\newcommand{\minrk}{\mathrm{minrk}}
\newcommand{\MAIS}{\mathrm{MAIS}}
\newcommand{\Crux}{\mathrm{Crux}}
\newcommand{\TCP}{\mathrm{TCP}}
\DeclareMathOperator{\rank}{rank}
\DeclareMathOperator{\dist}{dist}
\DeclareMathOperator{\poly}{poly}

% Notation specific to this written-proof manuscript.
\DeclareMathOperator{\Var}{Var}
\newcommand{\gaparea}{\mathcal A}

% Notation specific to this memo; macros.tex is an unchanged copy of the source-cover paper's.
\newcommand{\meanarea}{\overline{\gaparea}}
\newcommand{\limarea}{\gaparea^{*}}
\newcommand{\ZA}{Z^{\mathrm A}}
\newcommand{\ZB}{Z^{\mathrm B}}
\newcommand{\ZC}{Z^{\mathrm C}}
\newcommand{\vdc}{\mathrm{vdC}}
\setlist{itemsep=3pt,topsep=5pt}
\setlength{\emergencystretch}{2em}
\title{Frozen Gap Areas and the Sharpness of Source-Cover Certificates}
\author{Hans Brende \and Claude Opus 5.5 (Anthropic)}
\date{October 2026}
\begin{document}
\maketitle
\begin{abstract}
We continue the source-cover analysis of linear-time transposition of a $d\times d$
matrix with $b(d)=\max(1,\lceil\log_2d\rceil)$-bit entries on a multitape Turing
machine. For a fixed admissible return geometry, the source-cover paper of this series
bounds the correctness deficit below by the mean circular gap area of the complete
raw-source banks, frozen at the first output time of each column, minus an explicit
revelation cost. We use this frozen gap area itself as the certificate. Along every
joint profile of every near-full sequence, its limiting mean over legal phases and
columns is at most $\beta\tau_R/(M-R)$, where $\tau_R$ is the limiting hop tail after
the first $R$ core families (Theorem~G). This implies the source-cover paper's
necessary reciprocal-cover bound. We apply Theorem~G to two classes of geometries, with
no envelope on the sizes of the core families.
A \emph{funnel} of depth $m$ is a chain of time sets in which each set is fed by at
most one family, and only from the next set, while the last set is fed by none. In no
near-full sequence do infinitely many geometries have funnels of a fixed depth that contain
every read of a raw interval of fixed positive relative length, whatever the fixed number of
exceptional families (Theorem~A). A \emph{single-sweep terminal-landing} geometry has two
designated families, an older and a newer one, and its exceptional families are among
them. It reads the input in one nondecreasing sweep during the lower region of the older
one, keeps the input head at one coordinate during the lower regions of the core families
and of the newer one and at every column root, and lets none of these families act during
the sweep. No near-full sequence consists of such geometries (Theorem~B). Theorem~G cannot
exclude the boundary case of the source-cover paper's power envelopes for two exceptional
families. An explicit sequence of admissible geometries with two exceptional families and
hop counts at most $512L/i^2$ has mean frozen gap areas at most $10^7/R^2$ at every
modulus and hop tails at least $1/(32R)$, so the necessary condition of Theorem~G holds on
it whenever $3.2\cdot10^8(M-R)\le\beta R$, for instance at $M=2R+2$ for every
$\beta\ge10^9$ (Theorem~Z). These geometries carry a single input: they are not near-full,
and they say what the certificate can see, not that a transposer exists. All numbered
theorems, lemmas, corollaries and definitions are formalized in Lean; the comparison with
the reciprocal-cover bound is a remark with a proof sketch only, and a few elementary
consequences stated in prose are labelled as not formalized. These are certificates for
specified geometries, not a decision of the transposition question.

\end{abstract}
\begin{quote}\small
\textbf{Contributions.} Hans Brende chose the problem and directed the research.
Claude Opus 5.5 (Anthropic) developed the mathematics, the Lean formalization and the
text. The memo extends and reuses the Lean library of the source-cover paper by Hans
Brende and GPT-6 Astra (OpenAI). All numbered theorems, lemmas, corollaries and
definitions are formalized in Lean; \cref{rem:reciprocal-comparison} is a proof sketch, and the outlines of
the formal proofs are explanatory guides.
\end{quote}
\begingroup\small\tableofcontents\endgroup
\section{Introduction}\label{sec:introduction}

\paragraph{The question.}
Let $b(d)=\max(1,\lceil\log_2d\rceil)$ and $L=d^2b(d)$. Question~(T) asks whether one
deterministic multitape Turing machine, with a fixed finite number of tapes and a
fixed finite alphabet, transposes every $d\times d$ matrix with $b(d)$-bit entries,
given and returned in row-major order, in $O(L)$ steps.

The source-cover paper of this series~\cite{SourceCovers2026} fixes one complete
computation geometry and a nonempty set $\mathcal C$ of inputs that it transposes
correctly. It bounds the correctness deficit $\Delta=L-\log_2|\mathcal C|$ below by a
geometric credit minus the information revealed by a certificate
\cite[Theorem~3.5]{SourceCovers2026}. At each legal phase $\rho$ and column $j$ the
credit is at most the circular point-gap area $\gaparea(S_\rho(s_j))$ of the complete
raw-source set of an exact source bank, frozen at the first output time $s_j$ of the
column \cite[Theorems~2.7 and~2.9]{SourceCovers2026}. That paper bounds the frozen area
below by fractional covers, anchored approximation and shared support. It obtains a
uniform exclusion for one exceptional family, and further exclusions under
economical-cover or shared-support conditions. For two exceptional families it
identifies the inverse-square hop profiles as the nearest unresolved case
\cite[Section~9.1]{SourceCovers2026}.

\paragraph{Results.}
\begin{enumerate}
\item \emph{Theorem~G (\cref{thm:necessary-gap-area}).} We keep the frozen point area
itself as the certificate. Along every joint profile of every near-full sequence, and
for all integers $R\ge0$ and $M\ge\max(2,R+1)$,
\[
 \limarea_{R,M}\le\frac{\beta\tau_R}{M-R},
\]
where $\limarea_{R,M}$ is the limiting mean frozen point area over legal phases and
columns, with the first $R$ core families protected (\cref{def:gap-area-means}). The
limits are those of the reciprocal-cover bound \cite[Theorem~4.6]{SourceCovers2026},
taken in the same order. Since the frozen point area dominates the fractional-cover
credit \cite[Lemma~3.3]{SourceCovers2026}, Theorem~G implies that bound
(\cref{rem:reciprocal-comparison}).
\item \emph{Theorem~A (\cref{thm:funnels}).} A funnel is a property of the full family
system: time sets $D_0,\dots,D_m$ such that each $D_i$ with $i<m$ is fed by at most
one family, which moves into $D_i$ only times of $D_{i+1}$, while no family feeds
$D_m$ (\cref{def:funnel}). Every orbit then meets $D_0$ in at most $m+1$ times
(\cref{lem:funnel-area}). With Theorem~G, no near-full sequence has, from some index on
or along a subsequence, funnels of a fixed depth whose first set contains every time
that reads a raw interval of fixed positive relative length. The number of exceptional
families and the sizes of the core families are arbitrary.
\item \emph{Theorem~B (\cref{thm:terminal-landings}).} A single-sweep terminal-landing
geometry (\cref{def:terminal-landing}) has two designated families, an older family
$e$ and a newer family $g$, and its exceptional families are among them. Its raw path
is constant on the lower regions of all core families and of $g$ and at all column
roots, and nondecreasing on the lower region of $e$; no core family and not $g$ acts
on that lower region. No near-full sequence consists, from some index on or along a
subsequence, of such geometries. No size envelope is assumed.
\item \emph{Theorem~Z (\cref{thm:sharpness,cor:sharpness-limit}).} At moduli
comparable to $R$, Theorem~G cannot exclude the two-exception class with an
inverse-square hop envelope $|G_i|\le HL/i^2$, the boundary case of the power envelopes
excluded in \cite[Corollary~4.4]{SourceCovers2026}. We construct admissible geometries
$\Gamma_\nu$ with two exceptional families and hop counts $|G_i|\le512L/i^2$
(\cref{def:sharp-geometry}). In every hop order, at every protected prefix $J\ge R$ and
every modulus $M>R$, their mean frozen point areas over legal phases and columns are at
most $10^7/R^2$, and their hop tails after $R$ are at least $1/(32R)$. Hence
$\limarea_{R,M}\le\beta\tau_R/(M-R)$ along every joint profile whenever
$3.2\cdot10^8(M-R)\le\beta R$, for instance at $M=2R+2$ for every $\beta\ge10^9$.
The geometries carry one input; frozen point areas and hop profiles do not involve the
correctness sets. Two exceptional families thus suffice to meet the frozen gap-area
condition at these moduli.
\end{enumerate}

\paragraph{How to read the memo and its formal proofs.}
Every numbered theorem, lemma and definition is formalized in Lean. Each carries a
status line naming its Lean file and declarations, and each theorem and lemma is
followed by an outline of the formal proof. The Lean statements, including their
inherited hypotheses, are authoritative; the outlines are guides.
\Cref{rem:reciprocal-comparison}, the comparison with the reciprocal-cover bound, is
not formalized: it is a remark with a proof sketch only. A few elementary consequences
stated in prose, after Theorems~G and~A and in \cref{sec:terminal-landings,sec:sharpness},
are not formalized either, and are labelled so where they occur. Results of the
source-cover paper are cited by their numbers at the pinned revision; all of them are
Lean-verified there. The memo's library imports that paper's library unchanged
(\cref{sec:verification}).

\paragraph{Related work.}
The model, the exact source bank, cyclic freezing, the paid inequality, joint profiles
and the necessary hop-tail and reciprocal-cover bounds are those
of~\cite{SourceCovers2026}. The finite form of Theorem~G is the case of the formal paid
inequality of \cite[Theorem~3.5]{SourceCovers2026},
\lean{SourceCovers.paid_deficit_of_gap_credit}, whose credit is the frozen point area
itself. The new content of Theorem~G is its limit form and its use in
\cref{sec:funnels,sec:terminal-landings}. The counting there rests on the orbit
structure of \cite[Lemma~2.6]{SourceCovers2026} and on the interval bound recalled in
\cref{sec:preliminaries}. We claim no novelty for these elementary covering and
counting principles.

\paragraph{Organization.}
\Cref{sec:preliminaries} recalls the notation. \Cref{sec:gap-area} proves Theorem~G,
\cref{sec:funnels} Theorem~A and \cref{sec:terminal-landings} Theorem~B.
\Cref{sec:sharpness} constructs the sharpness geometries and proves Theorem~Z, and
\cref{sec:remains} states what remains. The appendices index the Lean declarations and give the build
and check commands.

\section{Notation from the source-cover paper}\label{sec:preliminaries}

We use the notation of~\cite{SourceCovers2026} and recall only what the statements
below need. All objects in this section are defined and Lean-verified there.

\paragraph{Geometries.}
Fix integers $d,b\ge1$ and put $L=d^2b$, $W=db$ and $\Omega=\{0,\dots,L-1\}$; for~(T),
$b=b(d)$. A fixed complete geometry consists of the input-head and work-head walks,
the halting time $T$ and the append times $t_0<\dots<t_{L-1}<T$
\cite[Section~2.1]{SourceCovers2026}. The \emph{raw path} $x$ is the input-head
position clamped to $\Omega$; it moves by at most one cell per step. In the Lean model
times are natural numbers, and statements about all times quantify over all of them.
The first append time of column $j$ is $s_j=t_{jW}$
\cite[Theorem~2.9]{SourceCovers2026}. An \emph{admissible geometry} with parameters
$d,b,c,\beta,N,q,a_0$ \cite[Definition~2.2]{SourceCovers2026} adds a nonempty
correctness set $\mathcal C$, the runtime bound $T\le cL$, whole-step record capacity
at most $\beta$, an $N$-class span partition of the full family system, an original
$\mathrm Q_2$ core $\mathcal K_0$, a disjoint set $\mathcal E$ of at most $q$
exceptional complete families, and the bound $a_0$ on the hop count of every other
family. Its correctness deficit is $\Delta=L-\log_2|\mathcal C|$
\cite[(2.2)]{SourceCovers2026}.

\paragraph{Families and shadows.}
A complete actual family $F$ has an owner $o_F$, an outermost hop $(p_F,U_F)$, a full
span $(p_F,U_F]$ and a hop count $|F|$ \cite[Definition~2.1]{SourceCovers2026}.
We call $(p_F,o_F]$ its \emph{lower region} and $(o_F,U_F]$ its \emph{active
interval}. Its clipped shadow $s_F$ is the identity off the active interval and moves
every time of it; by \cite[Definition~2.4 and (2.6)]{SourceCovers2026},
\begin{equation}\label{eq:strict-shadow}
 p_F<s_F(t)\le o_F<t\le U_F\qquad\text{whenever }o_F<t\le U_F;
\end{equation}
on its active interval $s_F$ is antitone. We write $F(t)=s_F(t)$ and say that $F$
\emph{acts} at $t$ if $F(t)\ne t$, that is, if $t$ lies in its active interval.
For a set $\mathcal S$ of families, $\mathcal O_{\mathcal S}(t)$ is the set of values
obtained from $t$ by finite compositions of their shadows, including the empty
composition \cite[Definition~2.4]{SourceCovers2026}; its elements are at most $t$.
A word $(F_1,\dots,F_n)$ of families is \emph{successful} from $t$ if $F_1$ acts at
$t$, $F_2$ acts at $F_1(t)$, and so on; its \emph{endpoint} is
$F_n(\cdots F_1(t)\cdots)$.

\paragraph{Cores, phases and banks.}
A set $\mathcal K$ of complete families is $\mathrm Q_2$ if its spans split into two
classes of pairwise disjoint spans. Joining two families of $\mathcal K$ when their
spans intersect gives a caterpillar forest, the \emph{overlap forest}; its
\emph{original nonleaves} are the vertices with at least two neighbours. The nonleaves
of each component form a path, so a nonleaf has at most two nonleaf neighbours, and a
leaf has at most one neighbour \cite[Lemma~2.6]{SourceCovers2026}. Fix $M\ge2$ and a
protected set of at most $R<M$ families of $\mathcal K$. For a residue
$\rho\in\{0,\dots,M-1\}$, the original nonleaves whose rank along their nonleaf path
is congruent to $\rho$ modulo $M$ are deleted, and the other families of $\mathcal K$
\emph{survive}; they form the retained core. A residue is a \emph{legal phase} if no
protected nonleaf has a rank congruent to it; there are at least $M-R$ legal phases.
Let $A_\rho(u)$ be the orbit of the retained core at $u$. Every successful word of
surviving families has at most $M+1$ letters, and $|A_\rho(u)|\le2M+3$
\cite[Lemma~2.6]{SourceCovers2026}. At most two families of $\mathcal K$ act at any one
time, since their spans contain that time and so lie in distinct classes. For an
exceptional set $\mathcal E$, the exact source bank has value set $O_\rho(t)$ at $t$,
the orbit of the retained core together with the exceptional families
\cite[Theorem~2.7]{SourceCovers2026}. The raw source set is $S_\rho(t)=x(O_\rho(t))$
\cite[Section~3.1]{SourceCovers2026}; it contains $x(t)$. In limit statements,
$\mathcal K$ consists of the first $J$ families of the original core in its supplied
order by nonincreasing hop count, the protected set consists of the first $R$, and the
forest, ranks and phases are rebuilt for every selection
\cite[Section~4.2]{SourceCovers2026}.

\paragraph{Gap area and the paid error.}
View $\Omega$ as the integer points of a circle of circumference $L$. For a nonempty
$S\subseteq\Omega$ whose consecutive distinct points have circular gaps
$g_1,\dots,g_v$, the point-gap area is \cite[(2.14)]{SourceCovers2026}
\[
 \gaparea(S)=\frac1{2L^2}\sum_{h=1}^vg_h^2\in\Bigl[0,\frac12\Bigr].
\]
We call $\gaparea(S_\rho(s_j))$ the \emph{frozen point area} of phase $\rho$ and
column $j$. One counting fact is used repeatedly. If a nonempty $S\subseteq\Omega$
lies in the union of at most $n$ closed integer intervals $[a_i,b_i]$ with
$\sum_i(b_i-a_i)\le\alpha L$, then
\begin{equation}\label{eq:interval-budget}
 \gaparea(S)\ge\frac{(1-\alpha)_+^2}{2n},\qquad (z)_+=\max(z,0).
\end{equation}
This is the exposed-gap step in the proof of \cite[Lemma~3.3]{SourceCovers2026}:
an interval contains the left endpoint of at most one gap that it does not contain,
the exposed gaps therefore have total length at least $(1-\alpha)_+L$, and the
Cauchy--Schwarz inequality applies to them. Its Lean form is
\lean{LeanMachine.IntervalGapArea.point_area_lower_of_budgets} in the shared library.
The paid error of \cite[(3.7)]{SourceCovers2026} is
\[
 \mathsf E_q(w;J,M,R,a,\tau)=\frac{2C_q(T-1)}{dL}+\frac{2(w-1)}L+\frac{W-b+1}{2L}
 +\frac\beta L\Bigl[J+q+\frac\tau{M-R}+\frac{(1+Na)C_q(T-1)}w\Bigr],
\]
with $C_q=C_0(1+3C_0)^q$ and $C_0=16+24M(M-1)$.

\paragraph{Near-full sequences and joint profiles.}
A near-full sequence consists of admissible geometries indexed by $\nu$, with common
$c,\beta,N,q$, such that $d_\nu\to\infty$, $a_{0,\nu}/L_\nu\to0$ and
$\Delta_\nu/L_\nu\to0$ \cite[Definition~4.1]{SourceCovers2026}. Each geometry comes
with an order $G_1,G_2,\dots$ of its original core by nonincreasing hop count.
A joint profile is a coordinatewise limit $(h_i,v_i)_{i\ge1}$ of the normalized
profiles $h_i^{(\nu)}=|G_i|/L_\nu$ and $v_i^{(\nu)}=\Var(x_\nu;[p_{G_i},U_{G_i}])/L_\nu$
along one strictly increasing subsequence; both profiles are zero after the last
family. Write
\[
 \tau_R=\sum_{i>R}h_i,\qquad \tau_{R,J}=\sum_{R<i\le J}h_i,\qquad V_K=\sum_{i>K}v_i.
\]
The limiting profiles are summable, so $\tau_R\to0$ as $R\to\infty$. Limits in $\nu$
are taken along the subsequence of the profile. A strictly increasing subsequence of a
near-full sequence is near-full, and every near-full sequence has a joint profile
\cite[Section~4.1]{SourceCovers2026}.

\section{Theorem G: the necessary frozen gap area}\label{sec:gap-area}

The paid inequality of \cite[Theorem~3.5]{SourceCovers2026} accepts any credit that is
at most the frozen point area at each legal phase and column. Here the credit is the
frozen point area itself, and the limits are taken as for the reciprocal-cover bound
\cite[Theorem~4.6]{SourceCovers2026}.

\begin{definition}[Frozen gap-area means]\label{def:gap-area-means}
\formal{GapAreaMemo/GapArea.lean}{SourceCovers.phaseGapArea,SourceCovers.phaseGapArea_bounds,SourceCovers.meanGapArea,SourceCovers.AdmissibleGeometry.prefixMeanGapArea,SourceCovers.NearFullSequence.JointProfile.prefixGapArea,SourceCovers.NearFullSequence.JointProfile.gapAreaLimit,SourceCovers.NearFullSequence.JointProfile.gapAreaLimit_bounds}
Fix a geometry, a $\mathrm Q_2$ selected core, an exceptional set, $M\ge2$ and a
protected set of at most $R<M$ selected families. For a residue $\rho<M$ and a column
$j<d$, the frozen point area $\gaparea(S_\rho(s_j))$ lies in $[0,1/2]$. Let $\meanarea$
be its arithmetic mean over the legal phases $\rho$ and all columns $j$. For an admissible geometry and an order of its original core by
nonincreasing hop count, $\meanarea_{J,R,M}$ denotes this mean when the first $J$
original core families are selected and the first $R$ protected. Along a joint
profile of a near-full sequence, $\meanarea_{\nu,J,R,M}$ is this mean for the
geometry of index $\nu$ on the subsequence of the profile, and
\begin{equation}\label{eq:gap-area-limit}
 \limarea_{R,M}=\limsup_{J\to\infty}\,\limsup_{\nu\to\infty}\,\meanarea_{\nu,J,R,M},
 \qquad \limarea_R=\limarea_{R,2R+2}.
\end{equation}
For $R<M$ both upper limits are real, and $0\le\limarea_{R,M}\le1/2$.
\end{definition}

\begin{theorem}[Theorem G: necessary frozen gap area]\label{thm:necessary-gap-area}
\formal{GapAreaMemo/GapArea.lean}{SourceCovers.gap_area_paid_deficit,SourceCovers.AdmissibleGeometry.prefix_gap_area_deficit,SourceCovers.NearFullSequence.JointProfile.prefix_gap_area_limit,SourceCovers.NearFullSequence.JointProfile.necessary_gap_area,SourceCovers.NearFullSequence.JointProfile.necessary_gap_area_default}
\begin{enumerate}[label=(\alph*)]
\item \emph{Finite form.} Take the data of \cite[Theorem~3.5]{SourceCovers2026}: a
fixed complete geometry with a nonempty correctness set, whole-step capacity at most
$\beta$ and an $N$-class span partition of its full family system; a $\mathrm Q_2$
selected core of at most $J$ families; at most $q$ exceptional families; $M\ge2$ and a
protected set of at most $R<M$ selected families; a bound $a$ on the hop count of
every family outside the selected core and the exceptions; and a real $\tau$ at least
the hop mass of the unprotected original nonleaves of the selected core. Then, for
every integer $w\ge1$,
\begin{equation}\label{eq:finite-gap-area}
 \frac\Delta L\ge\meanarea-\mathsf E_q(w;J,M,R,a,\tau).
\end{equation}
No runtime bound is used.
\item \emph{Prefixes.} For an admissible geometry with any order $G_1,G_2,\dots$ of its
original core by nonincreasing hop count, and integers $J\ge R\ge0$,
$M\ge\max(2,R+1)$, $w\ge1$ and $a\ge a_0$ with $a\ge cL/(J+1)$,
\[
 \frac\Delta L\ge\meanarea_{J,R,M}
 -\mathsf E_q\Bigl(w;J,M,R,a,\sum_{R<i\le J}|G_i|\Bigr).
\]
\item \emph{Dimension limit.} Along every joint profile of every near-full sequence,
for all integers $J\ge R\ge0$ and $M\ge\max(2,R+1)$ and every real $\lambda>0$,
\[
 \limsup_{\nu\to\infty}\meanarea_{\nu,J,R,M}
 \le\frac{\beta\tau_{R,J}}{M-R}+2\lambda+\frac{\beta NC_qc^2}{\lambda(J+1)}.
\]
\item \emph{The limit bound.} For every joint profile of every near-full sequence and
all integers $R\ge0$ and $M\ge\max(2,R+1)$,
\begin{equation}\label{eq:GA}
 \limarea_{R,M}\le\frac{\beta\tau_R}{M-R}.\tag{GA}
\end{equation}
In particular $\limarea_R\le\beta\tau_R/(R+2)$.
\end{enumerate}
In (c) and (d) the cutoff $R$ and the modulus $M$ are fixed, and the number $q$ of
exceptional families is any fixed integer.
\end{theorem}

\begin{outline}
(a) Apply \lean{SourceCovers.paid_deficit_of_gap_credit}, the formal paid inequality
of \cite[Theorem~3.5]{SourceCovers2026}, with the frozen point area
$\gaparea(S_\rho(s_j))$ as the credit of each legal phase $\rho$ and column $j$. Its
premise, that each credit is at most the frozen point area, holds with equality. The
mean credit is then $\meanarea$, and the paid error is unchanged.

(b) The first $J$ original families form a $\mathrm Q_2$ selection of at most $J$
families, and the hop mass of its unprotected nonleaves is at most
$\sum_{R<i\le J}|G_i|$. An omitted core family has at most $cL/(J+1)$ hops, by the
sorted order and the total bound $\sum_i|G_i|\le cL$; every other non-exceptional
family has at most $a_0$ hops. Hence (a) applies with this $a$, as in the proof of
\cite[Theorem~4.2]{SourceCovers2026}.

(c) For fixed $J,R,M,\lambda$, the error-limit lemma of the source-cover paper,
\lean{SourceCovers.NearFullSequence.JointProfile.prefix_error_limit}, supplies along
the subsequence widths $w_\nu\ge1$ and values $a_\nu$ admissible in (b), with widths
$\lceil\lambda L_\nu\rceil$ as in the proof of \cite[Theorem~4.2]{SourceCovers2026},
and upper bounds $F_\nu$ for the paid errors such that $\Delta_\nu/L_\nu+F_\nu$ tends
to the right side of~(c). By (b), $\meanarea_{\nu,J,R,M}\le\Delta_\nu/L_\nu+F_\nu$, and
the means lie in $[0,1/2]$; pass to the upper limit.

(d) Let $J\to\infty$ in (c): $\tau_{R,J}\to\tau_R$ by summability, and the last term
tends to zero. The finite upper limits lie in $[0,1/2]$, so the outer upper limit
preserves the inequality. Finally take $\lambda=\delta/2$ for every $\delta>0$ and
remove the additive $\delta$. For $M=2R+2$, $M-R=R+2$. As in
\cite[Theorem~4.6]{SourceCovers2026}, the dimension limit, the prefix limit and the
removal of the width error occur in that order.
\end{outline}

\paragraph{Why it matters.}
Theorem~G is the necessary condition that the paid inequality gives without choosing a
cover witness. For an impossibility proof it reduces the exclusion of a class of
admissible geometries to one lower bound: it suffices to show that every near-full
sequence in the class has a joint profile with $\limarea_{R,M}>\beta\tau_R/(M-R)$ for
some $R$ and $M$. Such a bound concerns the geometries alone, since frozen point areas
and hop profiles do not involve the correctness sets. Theorems~A and~B prove such
bounds for two classes. For a construction it is a requirement on the frozen banks.
For example, $\gaparea(S)\ge1/(2|S|)$ for every nonempty $S$ by the Cauchy--Schwarz
inequality. Along a joint profile of a near-full sequence with an inverse-square
limiting profile $h_i\le H/i^2$, for which $\tau_R\le H/R$ when $R\ge1$, the reciprocal
sizes $1/|S_\rho(s_j)|$ of the frozen raw sets therefore have limiting mean at most
$2\beta H/(R(R+2))$ at modulus $2R+2$ for every $R\ge1$, in the sense of
\eqref{eq:gap-area-limit}. This consequence is elementary and
is not formalized.

\begin{remark}[Comparison with the reciprocal-cover bound]\label{rem:reciprocal-comparison}
Proof sketch only: for every joint profile of every near-full sequence, and more
generally of every sequence of admissible geometries with $d\to\infty$ (the argument
does not use near-fullness, and $Q^*_{R,M}$ is defined along the profile as in
\cite[(4.7)]{SourceCovers2026}), and all integers $K,R\ge0$ and $M\ge\max(2,R+1)$,
\[
 \frac{(1-V_K)_+^2}{2(2K+1)}\,Q^*_{R,M}\le\limarea_{R,M},
\]
where $Q^*_{R,M}$ is the reciprocal-cover limit of \cite[(4.7)]{SourceCovers2026}.
Hence \eqref{eq:GA} implies the weighted reciprocal-cover bound
\cite[(4.8)]{SourceCovers2026}, and through it the necessary hop-tail bound
\cite[Theorem~4.2]{SourceCovers2026}. Indeed, for $J\ge\max(K,R)$ distinguish the first $K$
selected families. At every legal phase and column, \cite[Lemma~3.3]{SourceCovers2026}
gives
\[
 \gaparea(S_\rho(s_j))\ge\frac{(1-\alpha_\nu)_+^2}{2(2K+1)\kappa_{\nu,J,R,M,\rho,j}},
 \qquad \alpha_\nu=\sum_{K<i\le J}v_i^{(\nu)},
\]
with the fractional cover $\kappa$ of \cite[Definition~3.1]{SourceCovers2026}; its
formal form, \lean{SourceCovers.fractional_phase_gap}, is stated for the same point
area. Average over legal phases and columns. Along the profile
$\alpha_\nu\to\sum_{K<i\le J}v_i\le V_K$, and all averages lie in $[0,1]$, so the
dimension upper limit of $\meanarea_{\nu,J,R,M}$ is at least $(1-V_K)_+^2/[2(2K+1)]$
times that of the mean reciprocal cover. The prefix upper limit preserves this.
Theorem~G retains the full frozen point area, including what the fractional-cover
estimate gives up through the factors $(1-\alpha)_+^2$ and $2(2K+1)$.
\end{remark}

\section{Theorem A: funnels}\label{sec:funnels}

A funnel restricts how shadows of the full family system can enter the times at which
a raw interval is read. It involves no root, selection or phase.

\begin{definition}[Feeding and funnels]\label{def:funnel}
\formal{GapAreaMemo/Funnels.lean}{SourceCovers.ActualFamily.Feeds,SourceCovers.Funnel,SourceCovers.AdmissibleGeometry.HasFunnel}
Fix a geometry. A complete family $F$ \emph{feeds} a set $D$ of times if $F(t)\ne t$
and $F(t)\in D$ for some time $t$. For an integer $m\ge0$, a \emph{funnel of depth
$m$} consists of sets $D_0,\dots,D_m$ of times and complete families
$f_0,\dots,f_{m-1}$ such that
\begin{enumerate}[label=(F\arabic*)]
\item for each $i<m$, every complete family that feeds $D_i$ is $f_i$;
\item for each $i<m$ and every time $t$ with $f_i(t)\ne t$ and $f_i(t)\in D_i$, the
time $t$ lies in $D_{i+1}$;
\item no complete family feeds $D_m$.
\end{enumerate}
Here the complete families range over the whole family system: core, exceptional and
all other families. The formal definition indexes the sets and families by all
$i\ge0$, of which only $D_0,\dots,D_m$ and $f_0,\dots,f_{m-1}$ matter; it therefore
presupposes a nonempty family system. An admissible geometry \emph{has a funnel of
depth $m$ over a raw interval of relative length $\delta$} if there are such a funnel
and integers
$0\le z_0\le z_1\le L$ with $z_1-z_0\ge\delta L$ such that every time $s$ with
$z_0<x(s)<z_1$ lies in $D_0$.
\end{definition}

For $m=0$, conditions (F1) and (F2) are vacuous: a funnel of depth zero is a set $D_0$
of times into which no shadow moves any time, and every orbit meets it at most in its
root. For larger $m$, the set $D_i$ can be entered only through $f_i$ and only from
$D_{i+1}$, and the chain ends in a set that nothing enters.

\begin{lemma}[Lemma A: a funnel forces point area]\label{lem:funnel-area}
\formal{GapAreaMemo/Funnels.lean}{SourceCovers.Funnel.reach,SourceCovers.Funnel.reach_card,SourceCovers.Funnel.orbit_inter_subset,SourceCovers.phaseGapArea_ge_of_funnel,SourceCovers.AdmissibleGeometry.prefixMeanGapArea_ge_of_funnel}
Fix a geometry with a funnel of depth $m$.
\begin{enumerate}[label=(\alph*)]
\item For every set $\mathcal S$ of complete families and every time $s$, the orbit
$\mathcal O_{\mathcal S}(s)$ meets $D_0$ in at most $m+1$ times. More precisely,
$\mathcal O_{\mathcal S}(s)\cap D_0\subseteq P_m(s)$, where $P_0(s)=\{s\}$ and
$P_{n+1}(s)=\{s\}\cup f_{m-n-1}(P_n(s))$ for $n<m$, so that $|P_n(s)|\le n+1$.
\item Let $0\le z_0\le z_1\le L$ be integers such that every time $s$ with
$z_0<x(s)<z_1$ lies in $D_0$. Then for every $\mathrm Q_2$ selected core, every
exceptional set, every $M\ge2$, every residue $\rho<M$ and every column $j<d$,
\[
 \gaparea(S_\rho(s_j))\ge\frac{\bigl((z_1-z_0)/L\bigr)^2}{2(m+3)}.
\]
\item If an admissible geometry has a funnel of depth $m$ over a raw interval of
relative length $\delta\ge0$, then $\meanarea_{J,R,M}\ge\delta^2/(2(m+3))$ for every
hop order and all $J$, $R$ and $M\ge\max(2,R+1)$.
\end{enumerate}
\end{lemma}

\begin{outline}
(a) We show that every word over families of $\mathcal S$ whose endpoint from $s$ lies
in $D_{m-n}$ has its endpoint in $P_n(s)$, by induction on the word from its last
letter. The empty word ends at $s\in P_n(s)$, and a last letter that does not act
leaves the endpoint unchanged. Suppose the last letter $F$ moves the previous endpoint
$u$ into $D_{m-n}$. Then $F$ feeds $D_{m-n}$, which (F3) excludes if $n=0$. If
$n\ge1$, then $F=f_{m-n}$ by (F1), and $u\in D_{m-n+1}$ by (F2). By induction
$u\in P_{n-1}(s)$, so the endpoint lies in $f_{m-n}(P_{n-1}(s))\subseteq P_n(s)$.
Take $n=m$.

(b) By the exactness of the bank \cite[Theorem~2.7]{SourceCovers2026}, every raw value
of $S_\rho(s_j)$ strictly between $z_0$ and $z_1$ is the value of $x$ at an orbit
point in $D_0$; by (a) there are at most $m+1$ such values. Hence $S_\rho(s_j)$ lies in
$[0,z_0]$, $[z_1,L]$ and at most $m+1$ singletons: at most $m+3$ closed intervals of
total length $L-(z_1-z_0)$. Apply \eqref{eq:interval-budget} with
$\alpha=1-(z_1-z_0)/L$.

(c) The bound of (b) holds at every legal phase and column; there is at least one
legal phase. Average.
\end{outline}

\begin{theorem}[Theorem A: funnels are excluded]\label{thm:funnels}
\formal{GapAreaMemo/Funnels.lean}{SourceCovers.NearFullSequence.not_eventually_funnel,SourceCovers.NearFullSequence.not_frequently_funnel}
Fix an integer $m\ge0$ and a real $\delta>0$. In no near-full sequence do all
geometries from some index on have a funnel of depth $m$ over a raw interval of
relative length $\delta$, and in none do infinitely many of its geometries have one.
Funnels are as in \cref{def:funnel}; in particular the family system is nonempty.
The parameters $c,\beta,N,q$ of the sequence are arbitrary, and no envelope on the
sizes of the core families is assumed.
\end{theorem}

\begin{outline}
Choose a joint profile and put $B=\delta^2/(2(m+3))>0$. If the geometries have such
funnels from some index on, then so do those on the subsequence of the profile, and
\cref{lem:funnel-area}(c) with $R=0$ gives $\limarea_{0,M}\ge B$ for every $M\ge2$.
Theorem~G with $R=0$ gives $\limarea_{0,M}\le\beta\tau_0/M$. An $M\ge2$ with
$BM>\beta\tau_0$ gives a contradiction. If funnels occur in infinitely many geometries,
these geometries form a strictly increasing subsequence, which is again a near-full
sequence, and the first case applies to it.
\end{outline}

\paragraph{Why it matters.}
Theorem~A excludes a property of the full family system with no hypothesis on the
number of exceptional families or on the sizes of the core families. For an
impossibility proof, it shows that bounded-depth unique feeding of the reads of a
macroscopic raw interval already contradicts near-full correctness. For a construction
in the admissible class it is a requirement, which can be read off by following
feeders backwards. Let $X$ be the set of times that read the open interval
$(z_0,z_1)$, put $D_0=X$, and as long as $D_i$ is fed by exactly one family $f_i$, let
$D_{i+1}$ be the set of times that $f_i$ moves into $D_i$. If some $D_i$ with $i\le m$
is fed by no family, then $D_0,\dots,D_i$, followed by empty sets with arbitrary
families, form a funnel of depth $m$ over the interval; this needs a nonempty family
system, and an empty one is excluded directly by Theorem~G, since then every frozen raw
set is a single point, of area $1/2$, while $\tau_R=0$. So, in a near-full sequence,
from some index on and for every raw interval of length at least $\delta L$, this
backward chain either reaches, among $D_0,\dots,D_m$, a set fed by at least two
families, or has $D_0,\dots,D_m$ each fed by exactly one family. This reformulation is
elementary and is not formalized.

\section{Theorem B: single-sweep terminal landings}\label{sec:terminal-landings}

The class of this section has two designated families: an older family $e$, whose
lower region carries the only raw sweep, and a newer family $g$. We first count orbit
points of a $\mathrm Q_2$ selected core, then bound the frozen point area at one
selection and phase, and finally combine two selections with Theorem~G.

\subsection{Successful core words}

\begin{lemma}[Successful core words]\label{lem:core-words}
\formal{GapAreaMemo/CoreWords.lean}{SourceCovers.coreSuccess_prefix_or,SourceCovers.coreSuccess_internal_nonleaf,SourceCovers.core_root_count}
Let $\mathcal K$ be a $\mathrm Q_2$ set of complete families, and consider words of
families of $\mathcal K$.
\begin{enumerate}[label=(\alph*)]
\item Two words that are successful from the same time and have the same first letter
are comparable: one is a prefix of the other.
\item Every letter of a successful word other than its first and its last letter is
an original nonleaf of the overlap forest of $\mathcal K$.
\item Let $r$ be a time, $X$ a finite set of times and $\Phi\subseteq\mathcal K$.
Suppose that each $v\in X$ is the endpoint of a nonempty word $w_v$ that is
successful from $r$, and that the last letter of $w_v$ belongs to $\Phi$ whenever it is
a nonleaf. Then $|X|\le2|\Phi|+4$.
\end{enumerate}
Part (b) holds for every set of complete families.
\end{lemma}

\begin{outline}
(a) Induct on the words. After a nonidentity step of a family, the new time lies in
its span at or before its owner. That family no longer acts there, no other family of
its class acts there, and at most one family of the other class does
\cite[proof of Lemma~2.6]{SourceCovers2026}; the formal statement is
\lean{SourceCovers.core_unique_after}. So the second letters agree whenever both words
continue.

(b) Consecutive letters of a successful word have intersecting spans and strictly
decreasing owners, so the word is a simple path of the overlap forest
\cite[Lemma~2.6]{SourceCovers2026}. An internal letter has two distinct neighbours on
that path.

(c) Consider first the points whose words end in a nonleaf. Two such words with the
same first and last letters are comparable by (a), and the letters of a successful word
are distinct, so the words are equal. The first letter acts at $r$, and at most two
families act there; this gives at most $2|\Phi|$ points. For words that end in a leaf,
two words of length at least two with the same first letter are comparable by (a), and
a proper prefix would make its last letter, a leaf, an internal letter of the longer
word, contrary to (b). Recording the first letter and whether the length is one gives
at most four such points.
\end{outline}

\subsection{A newer family and a terminal older family}

\begin{lemma}[Adding a newer and a terminal older family]\label{lem:landing-orbits}
\formal{GapAreaMemo/LandingOrbits.lean}{SourceCovers.shadowOrbit_insert_newer,SourceCovers.shadowOrbit_insert_terminal,SourceCovers.raw_landing_antitone}
Let $\mathcal F$ be a set of complete families, $g$ and $e$ complete families, and $s$ a
time.
\begin{enumerate}[label=(\alph*)]
\item Every point of $\mathcal O_{\mathcal F\cup\{g\}}(s)$ lies in
$\mathcal O_{\mathcal F}(s)$, or in $\mathcal O_{\mathcal F}(g(u))$ for some
$u\in\mathcal O_{\mathcal F}(s)$ at which $g$ acts.
\item Suppose that no family of $\mathcal F$ acts at any time of the lower region
$(p_e,o_e]$. Then every point of $\mathcal O_{\mathcal F\cup\{e\}}(s)$ lies in
$\mathcal O_{\mathcal F}(s)$ or equals $e(u)$ for some $u\in\mathcal O_{\mathcal F}(s)$
at which $e$ acts.
\item If $x$ is nondecreasing on $(p_e,o_e]$, then $x(e(u'))\le x(e(u))$ whenever
$o_e<u\le u'\le U_e$.
\end{enumerate}
\end{lemma}

\begin{outline}
Parts (a) and (b) follow by induction over the letters of a word. (a) After a
nonidentity step of $g$ the time is at most $o_g$, and no later step increases it, so
$g$ never acts again. (b) After a nonidentity step of $e$ the time lies in
$(p_e,o_e]$ by \eqref{eq:strict-shadow}, where neither $e$ nor a family of
$\mathcal F$ acts; the point is final. (c) On its active interval the shadow of $e$ is
antitone with values in $(p_e,o_e]$; compose with the nondecreasing raw path.
\end{outline}

\subsection{Landings at one selection and phase}

\begin{definition}[Newer closure, landings and starts]\label{def:landings}
\formal{GapAreaMemo/LandingCounts.lean}{SourceCovers.newerRoots,SourceCovers.newerClosure,SourceCovers.landingSet,SourceCovers.IsLandingStep,SourceCovers.stepJump,SourceCovers.stepMax,SourceCovers.landingStarts}
Fix a $\mathrm Q_2$ selected core $\mathcal K$, $M\ge2$, a residue $\rho<M$, complete
families $g$ and $e$, and a time $s$. The \emph{newer roots} of $s$ are
$N_\rho(s)=\{s\}\cup\{g(u):u\in A_\rho(s),\ g(u)\ne u\}$, and its \emph{newer closure}
is $C_\rho(s)=\bigcup_{r\in N_\rho(s)}A_\rho(r)$. Since orbits are closed under
composition, this is the set $C_g(s)$ of \cite[Proposition~3.7(b)]{SourceCovers2026},
with $g$ in the role of the newer exceptional map. Its \emph{landing times} are
$\Lambda_\rho(s)=C_\rho(s)\cap(o_e,U_e]$. A \emph{landing step} is a pair $(u,F)$ of a
landing time $u$ and a surviving family $F$ that acts at $u$ with
$F(u)\in(o_e,U_e]$. Its \emph{jump} is $x(e(F(u)))-x(e(u))$, truncated at zero, and
$\Gamma_\rho(s)$ is the largest jump of a landing step, or zero if there is none.
The \emph{starts} $\Lambda^0_\rho(s)$ are the landing times $v$ for which there is no
landing time $u$ and surviving family $F$ acting at $u$ with $F(u)=v$.
\end{definition}

\begin{lemma}[Counts at one selection and phase]\label{lem:landing-counts}
\formal{GapAreaMemo/LandingCounts.lean}{SourceCovers.newerClosure_count,SourceCovers.newerClosure_count_span,SourceCovers.landingStarts_card,SourceCovers.landing_cover}
In the setting of \cref{def:landings}:
\begin{enumerate}[label=(\alph*)]
\item If $D$ is a set of times and $\Phi\subseteq\mathcal K$ contains every original
nonleaf whose lower region meets $D$, then $|C_\rho(s)\cap D|\le(2M+4)(2|\Phi|+5)$.
\item If $F\in\mathcal K$, $p_F<a$ and $v\le U_F$, then
$|C_\rho(s)\cap(a,v)|\le11(2M+4)$.
\item $|\Lambda^0_\rho(s)|\le5(2M+4)$.
\item If $x$ is nondecreasing on $(p_e,o_e]$, then every $v\in\Lambda_\rho(s)$ has a
start $s_0\in\Lambda^0_\rho(s)$ with
\[
 x(e(s_0))\le x(e(v))\le x(e(s_0))+(M+1)\,\Gamma_\rho(s).
\]
\end{enumerate}
\end{lemma}

\begin{outline}
(a) There are at most $1+|A_\rho(s)|\le2M+4$ newer roots. Fix one root $r$. Erasing
inactive letters \cite[Lemma~2.6]{SourceCovers2026}, every point of $A_\rho(r)$ other
than $r$ is the endpoint of a nonempty successful word of surviving families from $r$.
By \eqref{eq:strict-shadow} the endpoint lies in the lower region of the last letter,
so for an endpoint in $D$ a nonleaf last letter belongs to $\Phi$.
\Cref{lem:core-words}(c) leaves at most $2|\Phi|+4$ such points, and with $r$ at most
$2|\Phi|+5$ points per root.

(b) Take for $\Phi$ the family $F$ together with the nonleaves adjacent to $F$. A
nonleaf whose lower region meets $(a,v)\subseteq(p_F,U_F]$ is $F$ or has a span
meeting that of $F$. A nonleaf has at most two nonleaf neighbours and a leaf has at
most one neighbour, so $|\Phi|\le3$, and (a) gives $11(2M+4)$.

(c) Every start is a newer root, or lies in $A_\rho(r)$ for a newer root $r$ and
differs from it. In the second case it is the endpoint of a nonempty successful
survivor word from $r$. Its last step starts at a point of the newer closure that is
larger than the start and is not a landing time, since otherwise the start would be
reached by a landing step; so that point exceeds $U_e$. The span of the last letter
therefore contains the boundary between $U_e$ and $U_e+1$, which at most two families
of $\mathcal K$ do. The first letter acts at $r$, so it is one of at most two
families. Two such words with the same first and last letters are equal, by
\cref{lem:core-words}(a) and the distinctness of the letters of a successful word.
Each root thus contributes at most $1+4$ starts.

(d) Use strong induction on $U_e-v$. A landing time $v$ that is not a start equals
$F(u)$ for a landing time $u>v$ and a surviving family $F$; then $(u,F)$ is a landing
step, which raises the raw value by at most $\Gamma_\rho(s)$. The chain of steps from
a start $s_0$ to $v$ is a successful word of surviving families, so it has at most
$M+1$ letters \cite[Lemma~2.6]{SourceCovers2026}. The lower bound is
\cref{lem:landing-orbits}(c), since $v\le s_0$.
\end{outline}

\subsection{Single-sweep terminal-landing geometries}

\begin{definition}[Single-sweep terminal-landing geometry]\label{def:terminal-landing}
\formal{GapAreaMemo/LandingGaps.lean}{SourceCovers.AdmissibleGeometry.TerminalLanding}
Single-sweep terminal-landing data for an admissible geometry, with original core
$\mathcal K_0$ and exceptional set $\mathcal E$, consist of complete families $e$ and
$g$, the \emph{older} and the \emph{newer} family, and a natural number $x_0$, such that
\begin{enumerate}[label=(L\arabic*)]
\item $\mathcal E\subseteq\{e,g\}$;
\item $x(t)=x_0$ whenever $t$ lies in the lower region of a family of $\mathcal K_0$ or
of $g$;
\item $x(s_j)=x_0$ for every column $j<d$;
\item $x$ is nondecreasing on the lower region $(p_e,o_e]$ of $e$;
\item no family of $\mathcal K_0$, and not $g$, acts at any time of $(p_e,o_e]$.
\end{enumerate}
A geometry with such data is a \emph{single-sweep terminal-landing geometry}.
The families $e$ and $g$ need not be distinct, and need not both be exceptional.
Condition (L5) holds for $g$ whenever $o_e\le o_g$ (an elementary remark, not
formalized).
\end{definition}

Thus the input is read in one nondecreasing sweep during the lower region of $e$,
the input head is parked at $x_0$ at all other times that can contribute to a frozen
bank (\cref{lem:landing-gaps}(a)), and no family of a bank moves a landing of $e$
again.

\begin{lemma}[Finite gap bounds for terminal landings]\label{lem:landing-gaps}
\formal{GapAreaMemo/LandingGaps.lean}{SourceCovers.AdmissibleGeometry.TerminalLanding.raw_values,SourceCovers.AdmissibleGeometry.TerminalLanding.gap_of_step,SourceCovers.AdmissibleGeometry.TerminalLanding.gap_of_starts}
Let an admissible geometry carry single-sweep terminal-landing data $(e,g,x_0)$. Let
$\mathcal K\subseteq\mathcal K_0$ be a $\mathrm Q_2$ selected core, $M\ge2$, $\rho<M$ a
residue and $j<d$ a column. Landing sets refer to $\mathcal K$, $M$, $\rho$, $g$, $e$
and the root $s_j$, and $S_\rho(s_j)$ to the geometry's exceptional set.
\begin{enumerate}[label=(\alph*)]
\item Every element of $S_\rho(s_j)$ is $x_0$ or $x(e(v))$ for some
$v\in\Lambda_\rho(s_j)$.
\item Let $F\in\mathcal K$ act at a time $u<T$ with $o_e<u\le U_e$ and
$o_e<F(u)\le U_e$, and let $\gamma=x(e(F(u)))-x(e(u))$, truncated at zero. Then
\[
 \gaparea(S_\rho(s_j))\ge\frac{(\gamma/L)^2}{2\bigl(11(2M+4)+3\bigr)}.
\]
The family $F$ need not survive at $\rho$.
\item With $\Gamma=\Gamma_\rho(s_j)$,
\[
 \gaparea(S_\rho(s_j))\ge
 \frac{\bigl(1-5(2M+4)(M+1)\Gamma/L\bigr)_+^2}{2\bigl(5(2M+4)+1\bigr)}.
\]
\end{enumerate}
\end{lemma}

\begin{outline}
(a) By (L1), the orbit $O_\rho(s_j)$ of the retained core with the exceptions lies in
the orbit of the retained core together with $g$ and $e$. By (L5) and
\cref{lem:landing-orbits}(b), each of its points lies in the orbit of the retained
core with $g$, or equals $e(u)$ for a point $u$ of that orbit at which $e$ acts. A
point of the first kind is the root, where $x=x_0$ by (L3), or lies in the lower
region of a retained family or of $g$, where $x=x_0$ by (L2). For the second kind,
$u\in C_\rho(s_j)$ by \cref{lem:landing-orbits}(a), and $u\in(o_e,U_e]$ because $e$
acts at $u$.

(b) Since $F(u)<u$, \cref{lem:landing-orbits}(c) gives $x(e(u))\le x(e(F(u)))$. By (a)
and the same monotonicity, every raw value strictly between them is $x_0$ or
$x(e(v))$ for a landing time $v$ with $F(u)<v<u$. Since $F$ acts at $u$, $p_F<F(u)$
and $u\le U_F$, so \cref{lem:landing-counts}(b) leaves at most $1+11(2M+4)$ raw values
in that open interval. Cover the raw set by the two complementary closed intervals and
these singletons, and apply \eqref{eq:interval-budget}; the formal step is
\lean{SourceCovers.pointArea_ge_of_sparse}.

(c) By (a) and \cref{lem:landing-counts}(d), the raw set lies in $\{x_0\}$ and the
closed intervals $[x(e(s_0)),x(e(s_0))+(M+1)\Gamma]$ with $s_0\in\Lambda^0_\rho(s_j)$.
By \cref{lem:landing-counts}(c) these are at most $1+5(2M+4)$ intervals of total length
at most $5(2M+4)(M+1)\Gamma$. Apply \eqref{eq:interval-budget}.
\end{outline}

\subsection{Theorem B}

\begin{lemma}[Prefix bounds for terminal landings]\label{lem:landing-prefix}
\formal{GapAreaMemo/TerminalLandings.lean}{SourceCovers.AdmissibleGeometry.landingGamma,SourceCovers.AdmissibleGeometry.landingEnergy,SourceCovers.AdmissibleGeometry.landingEnergy_bounds,SourceCovers.AdmissibleGeometry.prefixMeanGapArea_ge_later,SourceCovers.AdmissibleGeometry.prefixMeanGapArea_ge_own}
Let an admissible geometry carry single-sweep terminal-landing data, and fix an order
of its original core by nonincreasing hop count. For $J_0\ge0$ and a column $j$, let $\Gamma_{J_0}(j)$ be the largest of the
jumps $\Gamma_\rho(s_j)$ over all selections of the first $J\le J_0$ original core
families with the first $R\le J$ protected and modulus $2R+2$, and over the legal
phases $\rho$ of each selection. The \emph{landing energy}
\[
 E(J_0)=\frac1d\sum_{j<d}\Bigl(\frac{\Gamma_{J_0}(j)}L\Bigr)^2
\]
lies in $[0,1]$. Put $f_R=20(R+2)(2R+3)$. For every $R\ge0$:
\begin{enumerate}[label=(\alph*)]
\item $\meanarea_{J,R,2R+2}\ge E(J_0)/(88R+182)$ for every $J\ge J_0$;
\item $\meanarea_{J,R,2R+2}\ge\bigl(1/2-2f_R^2E(J)\bigr)/\bigl(2(20R+41)\bigr)$ for every
$J\ge R$.
\end{enumerate}
\end{lemma}

\begin{outline}
At $M=2R+2$ the constants are $88R+182=2(11(2M+4)+3)$, $20R+41=5(2M+4)+1$ and
$f_R=5(2M+4)(M+1)$. Each jump is at most $L$, so $E(J_0)\le1$.

(a) Either $\Gamma_{J_0}(j)=0$, or it is the jump of an actual landing step $(u,F)$ at
a selection of the first $J'\le J_0$ families. Then $F$ also belongs to the first
$J\ge J_0$ families, and \cref{lem:landing-gaps}(b) at the selection $(J,R,2R+2)$ and
any of its legal phases gives
$\gaparea(S_\rho(s_j))\ge(\Gamma_{J_0}(j)/L)^2/(88R+182)$; the family need not survive
there. Average over legal phases and columns.

(b) For $J\ge R$, the selection $(J,R,2R+2)$ is among those that define
$\Gamma_J(j)$, so $\Gamma_\rho(s_j)\le\Gamma_J(j)$ at each of its legal phases. With
$y=\Gamma_J(j)/L$ and $\alpha=f_R\Gamma_\rho(s_j)/L\le f_Ry$,
\[
 (1-\alpha)_+^2\ge1-2\alpha\ge1-2f_Ry\ge\tfrac12-2f_R^2y^2,
\]
the last step because $2(f_Ry-\tfrac12)^2\ge0$. Apply \cref{lem:landing-gaps}(c) and
average over phases and columns.
\end{outline}

\begin{theorem}[Theorem B: single-sweep terminal landings are excluded]\label{thm:terminal-landings}
\formal{GapAreaMemo/TerminalLandings.lean}{SourceCovers.NearFullSequence.not_eventually_terminalLanding,SourceCovers.NearFullSequence.not_frequently_terminalLanding}
In no near-full sequence are all geometries from some index on single-sweep
terminal-landing geometries, and in none are infinitely many of its geometries such
geometries. The parameters $c,\beta,N,q$ of the sequence are arbitrary, and no envelope
on the sizes of the core families is assumed.
\end{theorem}

\begin{outline}
Fix a joint profile; along its subsequence the geometries eventually carry such data.

\emph{Step 1: the landing energies vanish.} For every $J_0$, the landing energy of the
geometry of index $\nu$ tends to zero along the subsequence. Otherwise it is at least
some $\eps>0$ infinitely often. By \cref{lem:landing-prefix}(a), for every $R$ and
every $J\ge\max(J_0,R)$ the inner upper limit of $\meanarea_{\nu,J,R,2R+2}$ is at least
$\eps/(88R+182)$, so $\limarea_R\ge\eps/(88R+182)$. Theorem~G gives
$\eps/(88R+182)\le\beta\tau_R/(R+2)$, and since $88R+182\le91(R+2)$, $\eps\le91\beta\tau_R$
for every $R$, contrary to $\tau_R\to0$.

\emph{Step 2: the selections themselves.} Fix $R$ and $J\ge R$. By Step~1, eventually
$2f_R^2E(J)\le1/4$, and \cref{lem:landing-prefix}(b) gives
$\meanarea_{\nu,J,R,2R+2}\ge1/(8(20R+41))$. Hence $\limarea_R\ge1/(8(20R+41))$, and
Theorem~G gives
\[
 \beta\tau_R\ge\frac{R+2}{8(20R+41)}\ge\frac1{164}\qquad(R\ge0),
\]
contrary to $\tau_R\to0$. If such geometries occur infinitely often, they form a
near-full subsequence, and the first case applies to it.
\end{outline}

\paragraph{Why it matters.}
The two-exception class with an inverse-square hop envelope is the nearest unresolved
case of \cite[Section~9.1]{SourceCovers2026}, and the necessary hop-tail bound
\cite[Theorem~4.2]{SourceCovers2026} is compatible with it. Theorem~B excludes, with no
envelope, the two-exception geometries in which the older family carries the only raw
sweep and receives terminal landings. A near-full sequence whose geometries have at
most two exceptional families must therefore, from some index on and for every choice
of $e$, $g$ and $x_0$ with $\mathcal E\subseteq\{e,g\}$, violate one of (L2)--(L5): it reads a coordinate
other than $x_0$ during the lower region of a core family or of $g$, or at a column
root; or its raw path decreases somewhere on the lower region of $e$; or a core family
or $g$ acts during that lower region. For an impossibility proof the argument shows
how Theorem~G is used at two scales: a fixed prefix bounds the landing jumps of every
larger selection (\cref{lem:landing-prefix}(a)), and at the selection itself few
starts with small jumps confine the frozen raw set to a few short intervals
(\cref{lem:landing-prefix}(b)).

\section{Sharpness for two exceptional families}\label{sec:sharpness}

Theorems~A and~B exclude two classes by proving lower bounds on frozen point areas. The
source-cover paper excludes every common power envelope $|G_i|\le HL/i^p$ with $p>q$
\cite[Corollary~4.4]{SourceCovers2026}; for two exceptional families the inverse-square
envelope $p=2$ is the boundary case. This section shows that the frozen gap-area condition
\eqref{eq:GA} cannot exclude that boundary case at moduli comparable to $R$. We construct
admissible geometries with two exceptional families and hop counts $|G_i|\le512L/i^2$
whose frozen point areas have mean at most $10^7/R^2$ over legal phases and columns, at
every modulus, and whose hop tails are at least $1/(32R)$. Hence \eqref{eq:GA} holds on
them whenever $3.2\cdot10^8(M-R)\le\beta R$, for instance at $M=2R+2$ for every
$\beta\ge10^9$ (\cref{thm:sharpness,cor:sharpness-limit}). The geometries carry a single
input. Frozen point areas and hop profiles do not involve the correctness sets, so this is
a statement about what the certificate can see; it constructs no transposer.

\paragraph{The idea.}
Two banks of zigzags, A and~B, carry the root of every column down through bank~A,
across the newer exceptional family $g$ into bank~B, and down through bank~B; the input
head does not move during them. Bank~A turns the column and the depth $p$ at which the
chain leaves it into an offset $o$, one of about $2^\ell$ values when the protected
zigzags have levels up to $\ell$. Bank~B carries every such offset to each of its
zigzags $\ZB_k$. A third bank, C, is swept once by the input head. The older exceptional
family $e$ sends the descent and the first rise of $\ZB_k$ into the second rise of the
bank-C zigzag $\ZC_k$, at relative levels $\vdc(k)\mp o/4$ around a centre given by the
base-$2$ van der Corput sequence. The landing levels of one column, about $4^\ell$ in
number, fill the unit interval up to $O(4^{-\ell})$ (\cref{lem:rows}); a single bank
would give only about $2^\ell$ of them. The chains of bank~C then copy every landing
level into every earlier zigzag of bank~C, so the frozen raw values fill the input range
except for one arc of length $O(L/2^\ell)$ (\cref{lem:dense-gaps}). Landings that fall
outside a second rise are parked elsewhere and only add points.

\begin{definition}[The dyadic sharpness geometries]\label{def:sharp-geometry}
\formal{GapAreaMemo/Z3DGeometry.lean}{SourceCovers.Z3D.Index.geometry,SourceCovers.Z3D.Index.runtime_le}
Fix an integer $\nu\ge3$ and put $d=2^{2^\nu}$, so that $b=b(d)=2^\nu$ and
$L=d^2b=2^{2^{\nu+1}+\nu}$. Each bank has zigzags with indices
$1\le k\le k_\nu=2^{\nu+1}-1$; the index $k$ has level $\ell(k)=\lfloor\log_2k\rfloor\le\nu$
and weight $w_k=4^{-\ell(k)-1}$. Banks A and~B have heights $w_kL$ and bank~C has
heights $8w_kL$; consecutive heights differ by a factor $1$ or~$4$. A zigzag of height
$h$ rises $h$ cells, descends $h$ cells and rises $h$ cells again; every leg moves one
cell per step, except that one rising leg may be slowed by a constant factor. Write
$\vdc$ for the base-$2$ van der Corput sequence, which reverses binary digits:
$\vdc\bigl(\sum_ik_i2^i\bigr)=\sum_ik_i2^{-i-1}$ for digits $k_i\in\{0,1\}$.
\begin{enumerate}[label=(\roman*)]
\item \emph{Bank C} consists of $\ZC_1,\dots,\ZC_{k_\nu}$ in this time order. $\ZC_1$
starts at time~$1$; for $k\ge2$, the first rise of $\ZC_k$ runs during the second rise
of $\ZC_{k-1}$, slowed by the factor $w_{k-1}/w_k$. Let $U_k$ be the trough of $\ZC_k$,
and let $U^*=U_{k_\nu}+8w_{k_\nu}L$ be the end of bank~C.
\item \emph{Bank B} consists of $\ZB_{k_\nu},\dots,\ZB_1$ in this time order, after
bank~C. The second rise of $\ZB_k$ runs during the first rise of $\ZB_{k-1}$, slowed by
the factor $w_{k-1}/w_k$; that of $\ZB_1$ is slowed by~$4$ and takes $L$ steps. Bank~B
ends at the trough $O_B$ of~$\ZB_1$.
\item \emph{Bank A} consists of zigzags $\ZA_k$ laid out like those of bank~B, after the
second rise of $\ZB_1$; it ends at the trough $O_A$ of $\ZA_1$. The outputs are appended
during the second rise of $\ZA_1$, at $t_z=O_A+1+z$ for $z<L$.
\item \emph{Heads.} One work head performs the zigzags of odd index of the three banks in
time order, another those of even index. Each zigzag starts at a fresh cell above all
cells its head has visited before. After
the halting time each of these heads performs padding zigzags; complete families are
formed from hops before the halting time only, so the padding adds no family. Two
further work heads perform one excursion each. The newer family $g$ rises at unit speed
from $O_B+1$ and descends at unit speed during bank~A, reaching level $i$ at time
$O_A-i$, for $i<h_g=2\sum_kw_kL$. The older family $e$ descends at unit speed during
bank~B, reaching level $l$ at time $O_B-l$, and rises earlier. Level $0$ departs at
time $0$; for $l\ge1$, the level that reaches the bank-B time at offset $0\le r<2w_kL$
from the start of $\ZB_k$ departs at the landing time
\[
 U_k+\psi_k\bigl(\lfloor4w_kL\,\vdc(k)\rfloor+w_kL-r\bigr),
\]
where $\psi_k(z)=2z$ for $0\le z\le4w_kL$ (a landing in the second rise of $\ZC_k$),
$\psi_k(z)=z$ for $z<0$ and $\psi_k(z)=4w_kL+z$ for $z>4w_kL$ (parking in the descents
of $\ZC_k$ and $\ZC_{k+1}$, or, for $k=k_\nu$, shortly after bank~C).
\item \emph{Input head.} It sweeps once across bank~C:
$x(t)=\min(\lfloor(t-1)/\sigma\rfloor,L-1)$ with $\sigma=\lfloor(U^*-1)/L\rfloor$.
\item \emph{Computation.} A single input and a one-point record alphabet.
\end{enumerate}
The original core consists of the upper and lower families of the $3k_\nu$ bank
zigzags; the exceptional families are $e$ and~$g$; and the bound on all other families
is $a_0=0$. We write $\Gamma_\nu$ for this geometry.
\end{definition}

The Lean namespace is \texttt{Z3D}, for three banks of zigzags with dyadic parameters. The
parameters are \lean{SourceCovers.Z3D.Index}, the heads
\lean{SourceCovers.Z3D.Index.heads}, the raw path and the append times
\lean{SourceCovers.Z3D.Index.raw} and \lean{SourceCovers.Z3D.Index.out}, the two
exceptional schedules \lean{SourceCovers.Z3D.Index.eSchedule} and
\lean{SourceCovers.Z3D.Index.gSchedule}, and the core
\lean{SourceCovers.Z3D.Index.core}. Each zigzag splices two scheduled excursions of the
shared library.

\begin{lemma}[Admissibility]\label{lem:sharp-admissible}
\formal{GapAreaMemo/Z3DFamilies.lean}{SourceCovers.Z3D.Index.family_mem,SourceCovers.Z3D.Index.core_disjoint,SourceCovers.Z3D.Index.exceptions_card,SourceCovers.Z3D.Index.outside_none}
For every $\beta\ge0$, $\Gamma_\nu$ is an admissible geometry with $c=128$, $N=3$,
$q=2$ and $a_0=0$. Every complete family is a core family, $e$ or~$g$; the core has a
two-class span partition and the full family system a three-class one.
\end{lemma}

\begin{outline}
The walk of each of the first two heads is a sequence of zigzags at fresh cells, and its
complete families are exactly the upper families of its zigzags (the returns of the
first rise and the descent, owned at the peak) and the lower ones (the returns of the
descent and the second rise, owned at the trough); each exceptional excursion is one
native family. Hence every family is a bank family or an exception. Upper families of
distinct bank zigzags have disjoint spans, and so do lower families; colouring the core
by upper and lower gives its two classes (\lean{SourceCovers.Z3D.Index.core_coloring}).
The spans of $e$ and $g$ are disjoint, and they form the third class
(\lean{SourceCovers.Z3D.Index.full_coloring}). The banks and the output phase take at
most $128L$ steps.
\end{outline}

\begin{lemma}[Inverse-square sizes and protection]\label{lem:sharp-sizes}
\formal{GapAreaMemo/Z3DProtected.lean}{SourceCovers.Z3D.Index.famA_card,SourceCovers.Z3D.Index.famB_card,SourceCovers.Z3D.Index.famC_card,SourceCovers.Z3D.Index.card_big_le,SourceCovers.Z3D.Index.hop_envelope,SourceCovers.Z3D.Index.famA_protected,SourceCovers.Z3D.Index.famB_protected,SourceCovers.Z3D.Index.famC_protected,SourceCovers.SortedCore.prefix_retained}
In $\Gamma_\nu$, the upper and lower families of $\ZA_k$ and of $\ZB_k$ have $w_kL$ hops
each, and those of $\ZC_k$ have $8w_kL$ hops. For every $\ell\ge0$, at most
$16\cdot2^\ell-6$ core families have at least $L\,4^{-\ell-1}$ hops. Hence, in every order
$G_1,G_2,\dots$ of the core by nonincreasing hop count:
\begin{enumerate}[label=(\alph*)]
\item $|G_i|\le512L/i^2$ for every $i$;
\item if $16\cdot2^\ell\le R$, the first $R$ families contain the upper and lower families
of every $\ZA_k$ and $\ZB_k$ with $\ell(k)\le\ell$ and of every $\ZC_k$ with
$\ell(k)\le\ell+1$; and for all $J\ge R$ and $M$, every family among the first $R$
survives at every legal phase.
\end{enumerate}
\end{lemma}

\begin{outline}
Each core family is the upper or lower family of one bank zigzag, and its hops are the
returns of one pair of legs, one per cell. The families with at least $L4^{-\ell-1}$ hops
belong to zigzags $\ZA_k$, $\ZB_k$ with $\ell(k)\le\ell$ and $\ZC_k$ with
$\ell(k)\le\ell+1$, at most $4(2^{\ell+1}-1)+2(2^{\ell+2}-1)$ families. (a) If
$L4^{-\ell-1}\le|G_i|<L4^{-\ell}$, then $G_1,\dots,G_i$ all have at least $L4^{-\ell-1}$
hops, so $i<16\cdot2^\ell$; and every family has at most $2L$ hops. (b) If a family with
at least $L4^{-\ell-1}$ hops stood at position $p>R$, the $p$ distinct families up to it
would all have that many hops, more than $16\cdot2^\ell$. A protected family survives at a
legal phase \cite[Lemma~2.6]{SourceCovers2026}.
\end{outline}

\begin{lemma}[Van der Corput rows with dyadic offsets]\label{lem:rows}
\formal{GapAreaMemo/Covering.lean}{SourceCovers.Covering.vdC,SourceCovers.Covering.vdC_hits,SourceCovers.Covering.w,SourceCovers.Covering.t,SourceCovers.Covering.t_sub_succ,SourceCovers.Covering.offsets,SourceCovers.Covering.offsets_dense,SourceCovers.Covering.row_cover_quarter,SourceCovers.Covering.rowPoints_dense_of_le}
For $p\ge1$ let $w_p$ be as in \cref{def:sharp-geometry}, put
$t_p=(2^{\ell(p)+1}-p)2^{-2\ell(p)-1}+2^{-\ell(p)-1}$, so that $t_1=1$ and
$t_p-t_{p+1}=2w_p$, and for $D\ge1$ and rational $\theta\in[0,1]$ let
$O_D(\theta)=\{t_p-\theta w_p,\ t_p-(2-\theta)w_p:1\le p\le D\}$. Let $n\ge6$ and
$2^n\le16D$. Then for every $k_0\ge0$ the points $\vdc(k)\pm o/4$ with $k_0\le k<k_0+2^n$
and $o\in O_D(\theta)$ meet every interval $[x,y]\subseteq[0,1]$ with rational endpoints
and $y-x\ge1152\cdot4^{-n}$.
\end{lemma}

\begin{outline}
Any $2^n$ consecutive values of $\vdc$ meet every dyadic interval of length $2^{-n}$,
hence every interval of length $2^{1-n}$ in $[0,1]$. The offsets lie in $(0,1]$, and above
every $r\in[t_D,1/2]$ some offset lies in $[r,r+8r^2]$, because consecutive values $t_p$
differ by $2w_p$ and $w_p\le r^2$ when $t_{p+1}\le r$. Given $[x,y]$, choose a centre
$\vdc(k)$ at a distance of order $2^{-n}$ from one end of the interval and an offset close
to four times that distance: the row point $\vdc(k)\pm o/4$ then lies in $[x,y]$. These
covering properties are elementary, and we claim no novelty for them.
\end{outline}

\begin{lemma}[Point area from dense values]\label{lem:dense-gaps}
\formal{GapAreaMemo/GapBound.lean}{SourceCovers.GapBound.window_cover,SourceCovers.GapBound.pointArea_le_of_dense'}
\begin{enumerate}[label=(\alph*)]
\item Let $\alpha_0,\dots,\alpha_n$ and $w'_0,\dots,w'_{n-1}$ be natural numbers with
$\alpha_s+w'_s\le\alpha_{s+1}$, and let $\Lambda\ge1$ be at least every $w'_s$ and every
$w'_s+w'_{s+1}$. If every window $[z,z+w'_s]\subseteq[\alpha_s,\alpha_{s+1}]$ meets a set
$S$ of natural numbers, then every window $[z,z+\Lambda]\subseteq[\alpha_0,\alpha_n]$
meets~$S$.
\item Let $S\subseteq\Omega$ be nonempty with $L-1\in S$, and let $\gamma\le X\le L$ be
integers such that every window $[u,u+\gamma]\subseteq[0,X]$ meets $S$. Then
\[
 \gaparea(S)\le\frac{(\gamma+1)L+(L-X+\gamma+1)^2}{2L^2}.
\]
\end{enumerate}
\end{lemma}

\begin{outline}
(a) Let $\alpha_s\le z<\alpha_{s+1}$. If $z+w'_s\le\alpha_{s+1}$, use the window
$[z,z+w'_s]$; otherwise use $[\alpha_{s+1},\alpha_{s+1}+w'_{s+1}]$, which lies in
$[z,z+\Lambda]$. (b) A gap starting at a point $u$ with $u+\gamma+1\le X$ has length at
most $\gamma+1$. The other gaps that do not wrap around lie in $[X-\gamma,L-1]$ and have
total length at most $L-X+\gamma+1$, and the wraparound gap ends at a point at
most~$\gamma$. Bound the sum of the squared gaps by the largest gap of each kind times the
total length of that kind.
\end{outline}

\begin{theorem}[Theorem Z: frozen gap areas of the sharpness geometries]\label{thm:sharpness}
\formal{GapAreaMemo/Z3DSharp.lean}{SourceCovers.Z3D.Index.phase_le,SourceCovers.Z3D.Index.mean_le,SourceCovers.Z3D.Index.gap_area_le,SourceCovers.Z3D.Index.hop_tail,SourceCovers.Z3D.Index.gap_area_condition,SourceCovers.Z3D.Index.gap_area_condition_default}
Let $\nu\ge3$ and $\beta\ge0$, and fix an order $G_1,G_2,\dots$ of the core of
$\Gamma_\nu$ by nonincreasing hop count.
\begin{enumerate}[label=(\alph*)]
\item If $1\le R\le2^\nu$, $M\ge R+1$ and $J\ge R$, then
$\meanarea_{J,R,M}\le10^7/R^2$.
\item If $R\ge1$, $4R\le2^\nu$ and $J\ge64R$, then $\sum_{R<i\le J}|G_i|\ge L/(32R)$.
\item Hence, if $R\ge1$, $4R\le2^\nu$, $J\ge64R$, $M\ge R+1$ and
$3.2\cdot10^8(M-R)\le\beta R$, then
\[
 \meanarea_{J,R,M}\le\frac{\beta}{M-R}\sum_{R<i\le J}\frac{|G_i|}L.
\]
In particular this holds for $M=2R+2$ whenever $\beta\ge10^9$.
\end{enumerate}
\end{theorem}

\begin{outline}
(a) For $R^2\le2\cdot10^7$ the bound is at least $1/2$. Otherwise choose $\ell\ge8$ with
$16\cdot2^\ell\le R<32\cdot2^\ell$. By \cref{lem:sharp-sizes}(b), at every legal phase the
orbit families contain the families of all bank zigzags of indices below $2^{\ell+1}$
and both exceptions. Fix a column $j\ge1$ and put $\theta=j/d$.
\begin{enumerate}[label=(\roman*)]
\item \emph{Bank A.} The root $s_j$ lies in the second rise of $\ZA_1$ at level
$\theta w_1L$. The lower family of $\ZA_1$ maps it into the descent, the upper family into
the first rise, which runs during the second rise of $\ZA_2$, and so on: $\ZA_p$ is
reached at the exact level $\theta w_pL$, in its descent and in its first rise
(\lean{SourceCovers.Z3D.Index.late_chain}).
\item \emph{Across $g$.} These times lie in the descent of $g$, which maps them into the
second rise of $\ZB_1$ at relative level $1-o$, with $o=t_p-\theta w_p$ from the
descent and $o=t_p-(2-\theta)w_p$ from the first rise, because
$2\sum_{i<p}w_i=1-t_p$.
\item \emph{Bank B and $e$.} The chain of bank~B reaches $\ZB_k$ at a level within $3$ of
$(1-o)w_kL$, the levels being rounded down. The family $e$ maps the descent time of
$\ZB_k$ into the second rise of $\ZC_k$ at a height $2\mu$ whose relative level
$(2\mu+1)/(8w_kL)$ lies within $1/(w_kL)$ of $\vdc(k)-o/4$, and the first-rise time to
one within $1/(w_kL)$ of $\vdc(k)+o/4$ (\lean{SourceCovers.Z3D.Index.landing}).
\item \emph{Density.} By \cref{lem:rows} with $k_0=2^\ell$, $n=\ell$ and $D=2^\ell$,
these relative levels, for $2^\ell\le k<2^{\ell+1}$, meet every window inside
$[2/(w_kL),1-2/(w_kL)]$ of relative length about $1152\cdot4^{-\ell}$
(\lean{SourceCovers.Z3D.Index.levels_window}).
\item \emph{Bank C.} The chain of bank~C carries a landing at height $\eta$ in $\ZC_k$ to
every $\ZC_i$ with $i\le k$, at the height $\eta'$ with
$(\eta'+1)/(8w_iL)=(\eta+1)/(8w_kL)$, in its second rise, its descent and its first rise
(\lean{SourceCovers.Z3D.Index.early_chain}). Hence every window of relative length about
$1152\cdot4^{-\ell}$ inside the first rise or the descent of each of
$\ZC_1,\dots,\ZC_{2^\ell}$ meets the orbit, and by \cref{lem:dense-gaps}(a) so does every
time window of length $4608L4^{-\ell}+O(4^\ell)$ that ends by the end of the first rise
of $\ZC_{2^\ell}$ (\lean{SourceCovers.Z3D.Index.time_cover}).
\item \emph{Raw values.} Since $x$ advances one cell every $\sigma\ge1$ steps, the frozen
raw set meets every window of that length inside $[0,X]$. Here
$L-X\le40L/2^\ell+O(L4^{-\ell})$, because the rest of bank~C lasts at most
$40\sum_{k>2^\ell}w_kL+16w_{2^\ell}L$ steps and $\sigma\ge2$; and $L-1=x(s_j)$ lies in the
set. \Cref{lem:dense-gaps}(b) gives $\gaparea(S_\rho(s_j))\le4010\cdot4^{-\ell}$
(\lean{SourceCovers.Z3D.Index.phase_bound}).
\end{enumerate}
None of this depends on $M$. Averaging over legal phases and columns, with the column
$j=0$ bounded by $1/2$, gives at most $4011\cdot4^{-\ell}<4.2\cdot10^6/R^2$. (b) Let
$2^r$ be the power of two with $2R<2^r\le4R$. The upper families of the zigzags $\ZA_k$
with $2^r\le k<2^{r+1}$ are $2^r$ distinct families of $L4^{-r-1}$ hops each. Fewer than
$16\cdot2^r\le J$ core families have that many hops, so all of them lie among the first $J$
(\cref{lem:sharp-sizes}); at most $R$ of them lie among the first $R$; hence
$\sum_{R<i\le J}|G_i|\ge(2^r-R)L4^{-r-1}\ge L/(32R)$. (c) Combine (a) and~(b):
$10^7/R^2\le\beta/(32R(M-R))$; for $M=2R+2$ use $M-R\le3R$.
\end{outline}

\begin{corollary}[Limit form]\label{cor:sharpness-limit}
\formal{GapAreaMemo/Z3DLimit.lean}{SourceCovers.Z3D.seq,SourceCovers.Z3D.limit_form,SourceCovers.Z3D.limit_GA,SourceCovers.Z3D.limit_GA_zero,SourceCovers.Z3D.limit_GA_default}
Consider the sequence $(\Gamma_{\nu+3})_{\nu\ge0}$ with any orders of its cores by
nonincreasing hop count, and any strictly increasing subsequence along which the
normalized hop profiles $|G_i|/L$ converge to limits $h_i$. Define $\limarea_{R,M}$ and
$\tau_R=\sum_{i>R}h_i$ as in \cref{def:gap-area-means} and \cref{sec:preliminaries}.
Then $\limarea_{R,M}\le10^7/R^2$ and $\tau_R\ge1/(32R)$ for all $R\ge1$ and $M\ge R+1$.
Hence
\[
 \limarea_{R,M}\le\frac{\beta\tau_R}{M-R}
\]
whenever $R\ge1$, $M\ge R+1$ and $3.2\cdot10^8(M-R)\le\beta R$, and whenever $R=0$ and
$2\le M\le\beta/16$. In particular \eqref{eq:GA} holds at $M=2R+2$ for every $R\ge0$
when $\beta\ge10^9$.
\end{corollary}

\begin{outline}
For fixed $J\ge R$, the geometries of the subsequence eventually satisfy $R\le2^\nu$, and
\cref{thm:sharpness}(a) bounds the inner upper limit; the outer one preserves the bound.
The finite profile masses are at most $128$, so the limit profile is nonnegative and
summable; \cref{thm:sharpness}(b) with $J=64R$ passes to the limit, and
$\tau_R\ge\sum_{R<i\le64R}h_i$. For $R=0$, $\limarea_{0,M}\le1/2$ and
$\tau_0\ge\tau_1\ge1/32$.
\end{outline}

\paragraph{Why it matters.}
For an impossibility proof, Theorem~Z closes one route: the frozen gap-area condition at
moduli up to $R+\beta R/(3.2\cdot10^8)$ cannot exclude the two-exception inverse-square
class. An exclusion of that class must use larger moduli, the remaining necessary
conditions of~\cite{SourceCovers2026}, or input specific to transposition
(\cref{prob:transposition-input}). For a construction, the geometries show how two
exceptional families can meet \eqref{eq:GA}: the older family lands the chains of the
protected banks on both sides of well-spread centres in a bank that the input head
sweeps, and the chains of that bank merge the landings of different zigzags.

The sequence of \cref{cor:sharpness-limit} belongs to the class of
\cref{prob:two-exception} for all $c\ge128$, $N\ge3$, $H\ge512$ and $\beta>0$, with
$a_0=0$, and \eqref{eq:GA} holds along every joint profile at every modulus allowed in
\cref{cor:sharpness-limit}. The argument of \cref{rem:reciprocal-comparison} does not
use near-fullness; applied to this sequence (proof sketch only), it shows that the
weighted reciprocal-cover bound \cite[(4.8)]{SourceCovers2026} holds at the same moduli.
By \cref{lem:funnel-area}(c), \cref{thm:sharpness} also implies that for large $\nu$ the
geometries have no funnel of a fixed depth over a raw interval of fixed positive
relative length; and they are not single-sweep terminal-landing geometries, since the
input head moves on the lower regions of the bank-C families. These two observations
are elementary and are not formalized. The constants are not optimized. Nothing here
concerns larger moduli, the necessary complete-observable movement and global
continuation density of \cite[Corollaries~4.8 and~5.3]{SourceCovers2026}, or the
existence of a transposer with such geometries.

\section{What remains}\label{sec:remains}

\paragraph{Scope.}
Theorems~G, A and~B are certificates for specified geometries. They assume the
admissible hypotheses of \cite[Definition~2.2]{SourceCovers2026}: one fixed full
geometry with a $\mathrm Q_2$ core, a bounded number of exceptional families, a bounded
full-span partition, and fixed runtime and capacity constants. These hypotheses are not
derived for an arbitrary linear-time transposer, and passing from an adaptive machine
to one geometry may lose entropy that must be paid for separately
\cite[Sections~1.1 and~9.3]{SourceCovers2026}. Nothing here decides~(T).

\paragraph{What a construction in the admissible class must do.}
For an oblivious transposer that is correct on every input, $\Delta=0$ on its single
geometry \cite[Section~1.1]{SourceCovers2026}. If its geometries are admissible with
fixed parameters and $a_0=o(L)$, they form a near-full sequence as $d\to\infty$; more
generally, every sequence of admissible geometries with $d\to\infty$, $a_0=o(L)$ and
$\Delta=o(L)$ is near-full. Such a sequence must satisfy
\eqref{eq:GA} at every protected scale, together with the necessary conditions of
\cite[Theorems~4.2 and~4.6 and Corollaries~4.8 and~5.3]{SourceCovers2026}, the last two
for two exceptional families. By Theorem~A, for every fixed depth $m$ and every fixed
$\delta>0$ its family systems eventually admit no funnel of depth $m$ containing all
reads of a raw interval of relative length $\delta$. By Theorem~B, its geometries
eventually admit no single-sweep terminal-landing data. These are necessary conditions,
not constructions of a transposer. By Theorem~Z, the condition \eqref{eq:GA} at moduli
$M\le R+\beta R/(3.2\cdot10^8)$ is compatible with two exceptional families and an
inverse-square hop envelope.

\paragraph{Open problems.}

\begin{problem}[Frozen gap areas for two exceptions]\label{prob:two-exception}
Fix $c,\beta>0$, $N\ge1$ and $H\ge0$. Consider sequences of admissible geometries with at most
two exceptional families and a common envelope $|G_i|\le HL/i^2$ in the supplied hop
order, with $d\to\infty$ and $a_0/L\to0$ but no condition on the correctness deficit.
Joint profiles and $\limarea_{R,M}$ are defined for them exactly as in
\cref{sec:preliminaries} and \cref{def:gap-area-means}. Does every such sequence have
a joint profile with $\limarea_{R,M}>\beta\tau_R/(M-R)$ for some $R$ and $M$? By
Theorem~G, a positive answer would exclude every near-full sequence in the class; a
negative answer would show that Theorem~G alone cannot exclude it. For $c\ge128$,
$N\ge3$ and $H\ge512$, the sequence of \cref{cor:sharpness-limit} shows that no pair
$(R,M)$ with $R\ge1$ and $3.2\cdot10^8(M-R)\le\beta R$, and no pair with $R=0$ and
$M\le\beta/16$, suffices; for $\beta\ge10^9$ this includes $M=2R+2$ for every $R$. The
question remains open for larger moduli relative to $R$ and for smaller constants. In
particular, for the sharpness geometries themselves it is open whether \eqref{eq:GA}
holds at every $R$ and $M$, and for which $\beta$.
\end{problem}

\begin{problem}[Transposition-specific input]\label{prob:transposition-input}
The certificates here use the transposition demands only through the cyclic-freezing
bound inside Theorem~G \cite[Theorem~2.9]{SourceCovers2026}. The curved-motion theorem
of the fourth paper, \cite[Theorem~3.12]{FourthPaper2026} (Theorem~Cq, proved there and
not formalized), states that in every macroscopic limit of a linear-time transposer the
work heads travel a total distance of at least half a block unit along motion that is
not countably-piecewise affine. Can such constraints on the work heads exclude, or
bound the frozen point areas of, admissible two-exception geometries that have no
funnels of bounded depth and are not single-sweep terminal-landing geometries? The
sharpness geometries of \cref{def:sharp-geometry} are a test case. Their work heads move
along affine legs, so in the macroscopic limit their paths are countably-piecewise affine.
Theorem~Cq does not apply to them as it stands: its hypothesis is an inequality that the
first paper's Theorem~M* proves for every limit of every machine
\cite[Section~3]{FourthPaper2026}, while $\Gamma_\nu$ carries a single input and is not
the run of a transposer. Bringing constraints of this kind into the admissible-geometry
framework is the open step. It is also not established here whether $\Gamma_\nu$ meets
the necessary complete-observable movement and global continuation density of
\cite[Corollaries~4.8 and~5.3]{SourceCovers2026}.
\end{problem}

Deciding (T) needs more than these certificates: for instance one fixed machine with a
proof that it transposes every input in $O(L)$ steps, or, as one possible route to the
negative answer, a reduction of every candidate machine to admissible geometries and
correctness sets with a positive net deficit, including the cost of selecting a geometry
in the adaptive case \cite[Section~9.3]{SourceCovers2026}.

\appendix
% Generated by tools/lean_index.py from the status lines; do not edit by hand.
\section{Index of Lean declarations}\label{sec:leanindex}

The status lines of the memo cite the following declarations, listed by Lean module
with their fully qualified names. Module files are relative to the memo's
\path{lean/} directory. The axiom checks in \path{lean/MemoAxioms.lean} are generated
from the same sources (\cref{sec:verification}).

\paragraph{\texttt{GapAreaMemo.CoreWords}}
{\small\raggedright\leanlist{SourceCovers.coreSuccess_internal_nonleaf,SourceCovers.coreSuccess_prefix_or,SourceCovers.core_root_count}.\par}

\paragraph{\texttt{GapAreaMemo.Covering}}
{\small\raggedright\leanlist{SourceCovers.Covering.offsets,SourceCovers.Covering.offsets_dense,SourceCovers.Covering.rowPoints_dense_of_le,SourceCovers.Covering.row_cover_quarter,SourceCovers.Covering.t,SourceCovers.Covering.t_sub_succ,SourceCovers.Covering.vdC,SourceCovers.Covering.vdC_hits,SourceCovers.Covering.w}.\par}

\paragraph{\texttt{GapAreaMemo.Funnels}}
{\small\raggedright\leanlist{SourceCovers.ActualFamily.Feeds,SourceCovers.AdmissibleGeometry.HasFunnel,SourceCovers.AdmissibleGeometry.prefixMeanGapArea_ge_of_funnel,SourceCovers.Funnel,SourceCovers.Funnel.orbit_inter_subset,SourceCovers.Funnel.reach,SourceCovers.Funnel.reach_card,SourceCovers.NearFullSequence.not_eventually_funnel,SourceCovers.NearFullSequence.not_frequently_funnel,SourceCovers.phaseGapArea_ge_of_funnel}.\par}

\paragraph{\texttt{GapAreaMemo.GapArea}}
{\small\raggedright\leanlist{SourceCovers.AdmissibleGeometry.prefixMeanGapArea,SourceCovers.AdmissibleGeometry.prefix_gap_area_deficit,SourceCovers.NearFullSequence.JointProfile.gapAreaLimit,SourceCovers.NearFullSequence.JointProfile.gapAreaLimit_bounds,SourceCovers.NearFullSequence.JointProfile.necessary_gap_area,SourceCovers.NearFullSequence.JointProfile.necessary_gap_area_default,SourceCovers.NearFullSequence.JointProfile.prefixGapArea,SourceCovers.NearFullSequence.JointProfile.prefix_gap_area_limit,SourceCovers.gap_area_paid_deficit,SourceCovers.meanGapArea,SourceCovers.phaseGapArea,SourceCovers.phaseGapArea_bounds}.\par}

\paragraph{\texttt{GapAreaMemo.GapBound}}
{\small\raggedright\leanlist{SourceCovers.GapBound.pointArea_le_of_dense',SourceCovers.GapBound.window_cover}.\par}

\paragraph{\texttt{GapAreaMemo.LandingCounts}}
{\small\raggedright\leanlist{SourceCovers.IsLandingStep,SourceCovers.landingSet,SourceCovers.landingStarts,SourceCovers.landingStarts_card,SourceCovers.landing_cover,SourceCovers.newerClosure,SourceCovers.newerClosure_count,SourceCovers.newerClosure_count_span,SourceCovers.newerRoots,SourceCovers.stepJump,SourceCovers.stepMax}.\par}

\paragraph{\texttt{GapAreaMemo.LandingGaps}}
{\small\raggedright\leanlist{SourceCovers.AdmissibleGeometry.TerminalLanding,SourceCovers.AdmissibleGeometry.TerminalLanding.gap_of_starts,SourceCovers.AdmissibleGeometry.TerminalLanding.gap_of_step,SourceCovers.AdmissibleGeometry.TerminalLanding.raw_values}.\par}

\paragraph{\texttt{GapAreaMemo.LandingOrbits}}
{\small\raggedright\leanlist{SourceCovers.raw_landing_antitone,SourceCovers.shadowOrbit_insert_newer,SourceCovers.shadowOrbit_insert_terminal}.\par}

\paragraph{\texttt{GapAreaMemo.TerminalLandings}}
{\small\raggedright\leanlist{SourceCovers.AdmissibleGeometry.landingEnergy,SourceCovers.AdmissibleGeometry.landingEnergy_bounds,SourceCovers.AdmissibleGeometry.landingGamma,SourceCovers.AdmissibleGeometry.prefixMeanGapArea_ge_later,SourceCovers.AdmissibleGeometry.prefixMeanGapArea_ge_own,SourceCovers.NearFullSequence.not_eventually_terminalLanding,SourceCovers.NearFullSequence.not_frequently_terminalLanding}.\par}

\paragraph{\texttt{GapAreaMemo.Z3DFamilies}}
{\small\raggedright\leanlist{SourceCovers.Z3D.Index.core_disjoint,SourceCovers.Z3D.Index.exceptions_card,SourceCovers.Z3D.Index.family_mem,SourceCovers.Z3D.Index.outside_none}.\par}

\paragraph{\texttt{GapAreaMemo.Z3DGeometry}}
{\small\raggedright\leanlist{SourceCovers.Z3D.Index.geometry,SourceCovers.Z3D.Index.runtime_le}.\par}

\paragraph{\texttt{GapAreaMemo.Z3DLimit}}
{\small\raggedright\leanlist{SourceCovers.Z3D.limit_GA,SourceCovers.Z3D.limit_GA_default,SourceCovers.Z3D.limit_GA_zero,SourceCovers.Z3D.limit_form,SourceCovers.Z3D.seq}.\par}

\paragraph{\texttt{GapAreaMemo.Z3DProtected}}
{\small\raggedright\leanlist{SourceCovers.SortedCore.prefix_retained,SourceCovers.Z3D.Index.card_big_le,SourceCovers.Z3D.Index.famA_card,SourceCovers.Z3D.Index.famA_protected,SourceCovers.Z3D.Index.famB_card,SourceCovers.Z3D.Index.famB_protected,SourceCovers.Z3D.Index.famC_card,SourceCovers.Z3D.Index.famC_protected,SourceCovers.Z3D.Index.hop_envelope}.\par}

\paragraph{\texttt{GapAreaMemo.Z3DSharp}}
{\small\raggedright\leanlist{SourceCovers.Z3D.Index.gap_area_condition,SourceCovers.Z3D.Index.gap_area_condition_default,SourceCovers.Z3D.Index.gap_area_le,SourceCovers.Z3D.Index.hop_tail,SourceCovers.Z3D.Index.mean_le,SourceCovers.Z3D.Index.phase_le}.\par}

\paragraph{Also named in the text}
{\small\raggedright From this memo's library: \leanlist{SourceCovers.Z3D.Index,SourceCovers.Z3D.Index.core,SourceCovers.Z3D.Index.core_coloring,SourceCovers.Z3D.Index.eSchedule,SourceCovers.Z3D.Index.early_chain,SourceCovers.Z3D.Index.full_coloring,SourceCovers.Z3D.Index.gSchedule,SourceCovers.Z3D.Index.heads,SourceCovers.Z3D.Index.landing,SourceCovers.Z3D.Index.late_chain,SourceCovers.Z3D.Index.levels_window,SourceCovers.Z3D.Index.out,SourceCovers.Z3D.Index.phase_bound,SourceCovers.Z3D.Index.raw,SourceCovers.Z3D.Index.time_cover,SourceCovers.pointArea_ge_of_sparse}.\par}
{\small\raggedright From the source-cover paper's library: \leanlist{SourceCovers.NearFullSequence.JointProfile.prefix_error_limit,SourceCovers.core_unique_after,SourceCovers.fractional_phase_gap,SourceCovers.paid_deficit_of_gap_credit}.\par}
{\small\raggedright From the shared library: \leanlist{LeanMachine.IntervalGapArea.point_area_lower_of_budgets}.\par}

\section{Building and checking the proofs}\label{sec:verification}

The memo's sources and PDF are licensed under CC BY 4.0, and its Lean library under
Apache 2.0. Readers need the whole repository, with sibling directories in place. The
memo's Lean package in \path{memo-gap-area/lean/} requires the source-cover paper's
package through the relative path \path{../../lean}. That package brings in its own
dependencies: the second paper's build adapter for the first paper's canonical sources,
the eighth paper's \path{boundedNesting} package and the shared \path{lean-machine}
library, and through them the second and fourth papers' libraries
(\path{transposeFollowup}, \path{motionDepthCodes}) and the second paper's pinned external
dependency, all at the repository's Lean and Mathlib v4.34.1 pins. Fetch or verify that
external dependency as described in the second paper's README. No Lean source of an
earlier paper is copied or modified; \path{paper/macros.tex} is an unchanged copy of the
source-cover paper's macros.

From the memo's \path{lean/} directory, run
\begin{verbatim}
lake exe cache get
lake build
lake env lean MemoAxioms.lean
\end{verbatim}
The cache command is needed only when the pinned dependency objects are not already
available. The cited declarations depend at most on \leanname{propext},
\leanname{Classical.choice} and \leanname{Quot.sound}. The memo's library uses no
\texttt{sorry}, no \texttt{native\_decide}, no new axiom and no heartbeat override.

From the memo directory, regenerate or check \cref{sec:leanindex} and
\path{lean/MemoAxioms.lean} with
\begin{verbatim}
python3 tools/lean_index.py
python3 tools/lean_index.py --check
\end{verbatim}
The script reads the status lines and the \texttt{\textbackslash lean} references of
the manuscript. It also checks that every numbered result has one status line followed
by the matching proof environment, that every declaration in a status line is declared
in the cited file, and that every declaration named in the text exists in one of the
libraries. It does not run Lean or check proofs.

Build the PDF from \path{paper/} with
\begin{verbatim}
latexmk -pdf -interaction=nonstopmode main.tex
\end{verbatim}

\begingroup\small\raggedright
\bibliographystyle{plainurl}
\bibliography{references}
\endgroup
\end{document}
